\section{Tool Calls as Executive Acts}
\label{sec:executive-acts}

A tool call is an executive act in a precise sense. The agent, holding delegated authority, dispatches that authority against a resource the agent does not own, producing an effect that, if wrong, must be undone. The structural anatomy is the same across domains.

\paragraph{The five-part structure.}
An executive act has five components. (1) An authorisation chain: who delegated authority to the acting party, under what scope, with what limitations. (2) A subject: the resource the act will affect. (3) An effect: a state transition or external communication produced by the act. (4) A claimed lifetime: the duration the act is intended to remain in force, after which it is expected to be discharged. (5) A reversal expectation: an implicit or explicit promise that the act can be undone if it should not have happened.

Tool calls instantiate each component. The authorisation chain is the credential and system prompt under which the agent runs, optionally enriched by a capability token or session policy. The subject is the tool server (a database, a filesystem path, an email recipient, a payments endpoint). The effect is whatever state the call mutates: a row inserted, a file written, a message dispatched. The lifetime is, in most current deployments, implicit and infinite: a write is permanent unless explicitly undone. The reversal expectation is implicit and often false: agent platforms typically log tool calls but rarely capture the inverse operation that would undo them.

\paragraph{The structural parallel.}
The same five components appear in domains that have been formally analysed in adjacent literature. An emergency response action in an endpoint-detection-and-response (EDR) system has an authorisation chain (the security policy plus an executor key), a subject (a process tree, a file, a network flow), an effect (quarantine, suspend, restrict egress), a claimed lifetime (a TTL after which the action is expected to be reviewed or withdrawn), and a reversal expectation (an inverse executor that can release the quarantine or resume the process). A temporally-bounded administrative order in a legal system has the same shape: a delegated authority, a subject, an effect, a sunset clause, and a reversal mechanism. An agent's call to delete a row from a customer database has the same structure.

The point of naming the parallel is not metaphor. The parallel is load-bearing because the formal grammar already exists for the EDR case and the temporally-bounded-order case, and the grammar is the same for the tool-call case. The bounded-executive-action construction (Section~\ref{sec:grammar}) captures all three under one typed envelope.

\paragraph{What the current deployment surface looks like.}
The Model Context Protocol~\cite{anthropicMCP2024} and analogous tool-calling protocols typically expose tools as a JSON-schema-described function call. The agent emits a function call; the runtime dispatches it to the tool server; the tool server executes and returns. Existing deployments may add an operator approval step, a denylist, or a per-tool rate limit. They do not, in general, require the tool's caller to construct the call only with a typed rollback witness. They do not enforce a positive TTL by construction. They do not produce a tamper-evident receipt that an external party can replay. The substrate is, in current deployments, structurally a passive dispatcher.

\paragraph{Why the executive-act framing matters for safety.}
Treating the tool call as an executive act exposes a class of guarantees the agent-as-policy framing obscures. A passive dispatcher cannot refuse a call on the basis of structural properties of the call itself; it can only refuse on properties of the agent that produced it (the agent's identity, rate limit, claimed prompt). A substrate that models the call as an executive act and demands its structural witnesses at admission can refuse calls that violate the grammar regardless of who produced them and regardless of whether the operator was manipulated into approving them.

The next section names the grammar explicitly.
